\documentclass[11pt]{article}

\usepackage[a4paper,margin=1in]{geometry}
\usepackage{amsmath,amssymb,amsthm,mathtools}
\usepackage{bm}
\usepackage{xcolor}
\usepackage{hyperref}
\usepackage{enumitem}
\usepackage{microtype}

\hypersetup{
  colorlinks=true,
  linkcolor=blue,
  urlcolor=blue,
  citecolor=blue
}

\newtheorem*{claim}{Claim}
\newtheorem*{lemma}{Lemma}
\newtheorem*{proposition}{Proposition}
\newtheorem*{remark}{Remark}
\newcommand{\R}{\mathbb{R}}
\newcommand{\E}{\mathbb{E}}
\newcommand{\norm}[1]{\left\lVert #1\right\rVert}
\newcommand{\inner}[2]{\left\langle #1,#2\right\rangle}

\title{Mathematical Theory for Deep Neural Networks\\HW3 Solution}
\author{Wells Guan\\University of Wisconsin--Madison}
\date{July 8, 2026}

\begin{document}
\maketitle

\section*{Problem 1. Weierstrass approximation and polynomial activations}

Since $\Omega\subset\R^d$ is closed and bounded, it is compact.

\subsection*{(a) Weierstrass Approximation Theorem}

The multivariable Weierstrass Approximation Theorem states that for every
$f\in C(\Omega)$ and every $\varepsilon>0$, there exists a polynomial
$p:\R^d\to\R$ such that
\[
  \norm{f-p}_{L^\infty(\Omega)}
  =\sup_{x\in\Omega}|f(x)-p(x)|<\varepsilon.
\]
Equivalently, if $\mathcal P(\R^d)$ denotes the set of all real polynomials in
$d$ variables, then
\[
  \overline{\mathcal P(\R^d)|_\Omega}^{\,\norm{\cdot}_{L^\infty(\Omega)}}
  =C(\Omega).
\]
Thus polynomials of arbitrarily high degree are dense in $C(\Omega)$.

\subsection*{(b) Networks with a polynomial activation}

Suppose that $\sigma$ is a polynomial of degree $r$. Consider a network with
$L$ hidden layers,
\[
  x^0=x,
  \qquad
  x^{\ell}=\sigma\bigl(W^{\ell-1}x^{\ell-1}+b^{\ell-1}\bigr),
  \quad \ell=1,\ldots,L,
\]
where $\sigma$ is applied componentwise, and with final affine output
\[
  f(x)=W^Lx^L+b^L.
\]
We show that every component of $x^\ell$ is a polynomial in $x$ of degree at
most $r^\ell$.

For $\ell=1$, each component of $W^0x+b^0$ is affine, and composition with the
polynomial $\sigma$ gives a polynomial of degree at most $r$. Suppose that the
claim holds at layer $\ell$. Then every component of
\[
  W^\ell x^\ell+b^\ell
\]
is a linear combination of polynomials of degree at most $r^\ell$, and hence
has degree at most $r^\ell$. Applying $\sigma$ increases the degree by at most
a factor $r$. Therefore every component of $x^{\ell+1}$ has degree at most
$r^{\ell+1}$. The claim follows by induction.

Consequently, every function in $\Sigma_{n_{1:L}}^\sigma$ is a polynomial of
degree at most $r^L$. If
\[
  \mathcal P_{r^L}(\Omega)
  :=\{p|_\Omega:p\text{ is a polynomial of degree at most }r^L\},
\]
then
\[
  \Sigma_{n_{1:L}}^\sigma\subset \mathcal P_{r^L}(\Omega).
\]
The space $\mathcal P_{r^L}(\Omega)$ is finite-dimensional and therefore closed
in $C(\Omega)$. When $\Omega$ is infinite, $C(\Omega)$ is infinite-dimensional,
so a finite-dimensional subspace cannot be dense in it. Hence
\[
  \overline{\Sigma_{n_{1:L}}^\sigma}^{\,\norm{\cdot}_{L^\infty(\Omega)}}
  \ne C(\Omega).
\]
Therefore a polynomial activation does not have the universal approximation
property for a fixed finite depth. The point is that width changes the number
of polynomial terms, but it does not remove the degree bound $r^L$.

\begin{remark}
If $\Omega$ is a finite set, then $C(\Omega)$ is finite-dimensional and the
answer may depend on the architecture and the locations of the points. The
non-density conclusion above is the standard statement for an infinite compact
domain, in particular for every compact domain with nonempty interior.
\end{remark}

\section*{Problem 2. Expected error of an empirical average}

Let
\[
  \Omega=[0,1],
  \qquad
  g_i(x):=g(x,w_i),
\]
where $w_1,\ldots,w_n$ are independent random variables with common density
$\lambda$. We assume that
\[
  \E_w\bigl[\norm{g(\cdot,w)}_{L^2(\Omega)}^2\bigr]<\infty,
\]
so that the following integrals and expectations are well defined.

\subsection*{(a) Formulas from the definition of expectation}

For every $x\in\Omega$, the mean function is
\[
  \mu(x)
  =\E_w[g(x,w)]
  =\int_{\R}g(x,t)\lambda(t)\,dt.
\]
The joint density of $(w_1,\ldots,w_n)$ is
\[
  \lambda_n(t_1,\ldots,t_n)
  =\prod_{i=1}^n\lambda(t_i).
\]
Therefore, by the definition of expectation,
\begin{align*}
  &\E_n\left[
    \norm{\mu-{1\over n}\sum_{i=1}^n g_i}_{L^2(\Omega)}^2
  \right] \\
  &\quad=
  \int_{\R^n}
  \norm{
    \mu(\cdot)-{1\over n}\sum_{i=1}^n g(\cdot,t_i)
  }_{L^2(\Omega)}^2
  \prod_{i=1}^n\lambda(t_i)\,dt_1\cdots dt_n \\
  &\quad=
  \int_{\R^n}\int_0^1
  \left|
    \mu(x)-{1\over n}\sum_{i=1}^n g(x,t_i)
  \right|^2
  dx\,
  \prod_{i=1}^n\lambda(t_i)\,dt_1\cdots dt_n.
\end{align*}
This is the required integral representation of the expected squared
$L^2$ error.

\subsection*{(b) Calculation of the expectation}

Set
\[
  Y_i:=g_i-\mu.
\]
Then $Y_1,\ldots,Y_n$ are independent, identically distributed random elements
of $L^2(\Omega)$ and
\[
  \E_n[Y_i(x)]=0
\]
for almost every $x\in\Omega$. Since
\[
  \mu-{1\over n}\sum_{i=1}^n g_i
  =-{1\over n}\sum_{i=1}^nY_i,
\]
we have
\begin{align*}
  \E_n\left[
    \norm{\mu-{1\over n}\sum_{i=1}^n g_i}_{L^2(\Omega)}^2
  \right]
  &= {1\over n^2}
  \E_n\left[
    \norm{\sum_{i=1}^nY_i}_{L^2(\Omega)}^2
  \right] \\
  &= {1\over n^2}
  \sum_{i,j=1}^n
  \E_n\bigl[\inner{Y_i}{Y_j}_{L^2(\Omega)}\bigr].
\end{align*}
For $i\ne j$, independence and the zero-mean property give
\begin{align*}
  \E_n\bigl[\inner{Y_i}{Y_j}_{L^2(\Omega)}\bigr]
  &=\int_0^1\E_n[Y_i(x)Y_j(x)]\,dx\\
  &=\int_0^1\E_n[Y_i(x)]\E_n[Y_j(x)]\,dx\\
  &=0.
\end{align*}
All diagonal terms are equal. Hence
\begin{align*}
  \E_n\left[
    \norm{\mu-{1\over n}\sum_{i=1}^n g_i}_{L^2(\Omega)}^2
  \right]
  &={1\over n}\E_w\bigl[
    \norm{g(\cdot,w)-\mu}_{L^2(\Omega)}^2
  \bigr].
  \tag{1}
\end{align*}
It remains to simplify the variance term. Expanding the square,
\begin{align*}
  \E_w\bigl[\norm{g(\cdot,w)-\mu}_{L^2(\Omega)}^2\bigr]
  &=\E_w\bigl[\norm{g(\cdot,w)}_{L^2(\Omega)}^2\bigr]
    -2\E_w\bigl[\inner{g(\cdot,w)}{\mu}_{L^2(\Omega)}\bigr]
    +\norm{\mu}_{L^2(\Omega)}^2.
\end{align*}
Since $\E_w[g(\cdot,w)]=\mu$,
\[
  \E_w\bigl[\inner{g(\cdot,w)}{\mu}_{L^2(\Omega)}\bigr]
  =\inner{\mu}{\mu}_{L^2(\Omega)}
  =\norm{\mu}_{L^2(\Omega)}^2.
\]
Therefore
\[
  \E_w\bigl[\norm{g(\cdot,w)-\mu}_{L^2(\Omega)}^2\bigr]
  =\E_w\bigl[\norm{g(\cdot,w)}_{L^2(\Omega)}^2\bigr]
   -\norm{\mu}_{L^2(\Omega)}^2.
\]
Substituting this identity into (1) yields
\[
  \boxed{
  \E_n\left[
    \norm{\mu-{1\over n}\sum_{i=1}^n g_i}_{L^2(\Omega)}^2
  \right]
  ={1\over n}\left(
    \E_w\bigl[\norm{g(\cdot,w)}_{L^2(\Omega)}^2\bigr]
    -\norm{\mu}_{L^2(\Omega)}^2
  \right).}
\]
Since $\norm{\mu}_{L^2(\Omega)}^2\ge0$, we also have the general upper bound
\[
  \E_n\left[
    \norm{\mu-{1\over n}\sum_{i=1}^n g_i}_{L^2(\Omega)}^2
  \right]
  \le {1\over n}
  \E_w\bigl[\norm{g(\cdot,w)}_{L^2(\Omega)}^2\bigr].
\]
Under the normalization
\[
  \E_w\bigl[\norm{g(\cdot,w)}_{L^2(\Omega)}^2\bigr]\le1,
\]
which is implicit in the last inequality stated in the problem, this becomes
\[
  \boxed{
  \E_n\left[
    \norm{\mu-{1\over n}\sum_{i=1}^n g_i}_{L^2(\Omega)}^2
  \right]\le {1\over n}.}
\]
Thus the mean squared $L^2$ error of the empirical average decays at the Monte
Carlo rate $n^{-1}$.

\end{document}
